\chapter{Mẹo thực hành và tài nguyên}

\begin{quote}\small

\textbf{Tóm tắt.} Chương 7 cung cấp một hướng dẫn thực hành về việc sử dụng AI tạo sinh trong môi trường học thuật, đưa ra các chiến lược tích hợp AI vào nghiên cứu, giảng dạy và các quy trình hành chính. Chương này phác thảo những nguyên tắc nền tảng để sử dụng AI hiệu quả, bao gồm kỹ thuật viết prompt, đánh giá có phê phán các kết quả đầu ra, và lựa chọn công cụ dựa trên bối cảnh và lĩnh vực chuyên môn. Chương nhấn mạnh tầm quan trọng của việc phát triển năng lực hiểu biết về AI (AI literacy), coi AI là một cộng tác viên chứ không phải là vật thay thế, đồng thời hướng dẫn cách xử lý những hạn chế thường gặp, chẳng hạn như hiện tượng ảo giác, lỗi trích dẫn và phân tích nông cạn. Bằng cách trau dồi các kỹ năng AI có khả năng thích ứng và vận dụng chúng một cách thận trọng, giới học thuật có thể nâng cao năng suất, duy trì liêm chính học thuật và luôn đi trước trong một bối cảnh công nghệ không ngừng biến đổi.

\textbf{Từ khóa:} Hiểu biết về AI, kỹ thuật viết prompt, năng suất học thuật, tích hợp AI, lỗi trích dẫn, tự động hóa quy trình làm việc, công cụ nghiên cứu, giảng dạy với AI, liêm chính học thuật, quản trị học thuật

\end{quote}

AI tạo sinh đã mang lại một bước chuyển mang tính biến đổi trong cách giới học thuật tiếp cận nghiên cứu, viết lách, giảng dạy và các công việc hành chính (Dwivedi et al., 2023). Trong khi các chương trước đã khám phá cách AI có thể hỗ trợ việc soạn thảo và trau chuốt các công trình học thuật, chương này mở rộng sang các ứng dụng rộng hơn của AI. Ngoài việc viết, AI có thể tinh giản quy trình làm việc, hỗ trợ quản lý nghiên cứu, cải thiện chiến lược giảng dạy và nâng cao hiệu quả trong các nhiệm vụ hành chính. Tuy nhiên, để khai thác trọn vẹn tiềm năng của AI, các học giả phải phát triển năng lực sử dụng những công cụ này, đánh giá một cách phê phán các kết quả đầu ra của chúng và tích hợp chúng một cách thận trọng vào công việc của mình (Bockting et al., 2023).

Chương này cung cấp một hướng dẫn thực tiễn về cách sử dụng AI tạo sinh hiệu quả trong môi trường học thuật. Chương đề cập các chiến lược lựa chọn và đánh giá công cụ AI, phát triển hiểu biết về AI, và tích hợp AI vào các công việc học thuật hằng ngày. Ngoài ra, chúng tôi thảo luận về những hạn chế và thách thức đi kèm với việc sử dụng AI, bao gồm các lỗi thường gặp (Southworth et al., 2023), vốn cũng sẽ được bàn thêm trong Chương 9. Vì công nghệ AI phát triển nhanh chóng, trọng tâm ở đây không phải là các công cụ cụ thể mà là những nguyên tắc chung sẽ vẫn hữu ích bất kể công nghệ tiến bộ như thế nào.

Hiểu cách sử dụng AI hiệu quả đòi hỏi nhiều hơn là chỉ quen thuộc với một mô hình hay nền tảng cụ thể. Cũng như những người dùng internet thuở đầu phải làm quen với các công cụ tìm kiếm không ngừng thay đổi, giới học thuật ngày nay phải thích ứng với một bối cảnh AI luôn biến động. Chìa khóa cho thành công lâu dài nằm ở việc phát triển các kỹ năng AI có khả năng thích ứng, đánh giá một cách phê phán nội dung do AI tạo ra, và sử dụng các công cụ này như những cộng sự chứ không phải sự thay thế cho chuyên môn học thuật.

\section{Bắt đầu với AI}

Chúng ta đã đề cập nhiều khía cạnh của AI tạo sinh (generative AI) trong các chương trước, đặc biệt liên quan đến việc soạn thảo các công trình học thuật. Chương này tiếp cận theo hướng rộng hơn, đưa ra những hướng dẫn chung về cách tích hợp hiệu quả AI tạo sinh vào các công việc học thuật ngoài viết lách. Chúng ta sẽ khám phá tiềm năng của nó trong nghiên cứu, giảng dạy và công tác hành chính, đồng thời bàn về cách đánh giá và lựa chọn những công cụ AI tốt nhất trong thời đại công nghệ phát triển nhanh chóng. Khi các mô hình và nền tảng mới liên tục xuất hiện, biết cách định hướng trong bối cảnh không ngừng biến đổi này là điều thiết yếu đối với những học giả muốn khai thác tối đa các công cụ ấy (Southworth et al., 2023).

Bước đầu tiên để sử dụng AI hiệu quả là làm quen với các công cụ AI tạo sinh đa dụng. Tuy nhiên, trước khi đi sâu vào từng nền tảng cụ thể, cần thừa nhận một hạn chế quan trọng: mọi điều chúng tôi viết ở đây đều có thể thay đổi nhanh chóng. Cũng như một cuốn sách về internet viết vào đầu thập niên 1990 có thể bàn về các công cụ tìm kiếm như Lycos hay WebCrawler, thì đến đầu những năm 2000, một cuốn hướng dẫn như vậy đã lỗi thời, vì Google và các công cụ tìm kiếm tiên tiến khác đã trở nên thống trị. Dù vậy, các nguyên tắc cốt lõi về cách sử dụng AI hiệu quả về cơ bản vẫn sẽ không đổi, ngay cả khi công nghệ tiếp tục tiến hóa. Cách chúng ta tương tác với AI tạo sinh, thông qua các prompt (câu lệnh) và việc tinh chỉnh lặp lại, nhiều khả năng sẽ vẫn giữ vai trò trung tâm, ngay cả khi bản thân các mô hình ngày càng được cải thiện.

Một trong những kỹ năng quan trọng nhất khi làm việc với AI là kỹ thuật viết prompt (prompt engineering). Dù đã đề cập vấn đề này ở Chương 3, vẫn cần nhấn mạnh lại: chất lượng đầu ra do AI tạo ra chịu ảnh hưởng trực tiếp từ độ chính xác của prompt của bạn. Nếu không hiểu AI tạo sinh có thể và không thể làm gì, bạn nhiều khả năng sẽ gặp khó khăn khi dùng nó, lãng phí thời gian hoặc nhận được kết quả kém (Mollick \& Mollick, 2023). Để tránh điều này, chúng tôi đặt ra một vài quy tắc (hay nguyên tắc) cơ bản giúp bạn sử dụng AI tạo sinh hiệu quả hơn trong công việc học thuật (hoặc bất kỳ công việc nào) của mình.

Quy tắc đầu tiên để sử dụng AI hiệu quả là phải cụ thể. Dù bạn đang dùng ChatGPT (có lẽ là ``Lycos'' của thời đại chúng ta), Claude (là ``WebCrawler''), hay DeepSeek, bạn đều phải truyền đạt rõ ràng điều mình cần. Bạn đang tìm kiếm các nghiên cứu liên quan? Động não ý tưởng? Soạn thảo văn bản? Đánh giá bài làm của sinh viên? Mỗi nhiệm vụ đòi hỏi một cách tiếp cận khác nhau, và các công cụ AI không thể suy đoán chính xác nhu cầu của bạn trừ khi bạn xác định chúng một cách tường minh (Southworth et al., 2023).

Điều này dẫn đến nguyên tắc cơ bản thứ hai: AI không thực sự hiểu bạn hay nội dung của bạn. Các mô hình AI tạo sinh là mô hình ngôn ngữ, không phải công cụ suy luận. Chúng vận hành dựa trên xác suất thống kê, dự đoán từ hoặc cụm từ tiếp theo có khả năng cao nhất để đáp lại đầu vào của bạn. Điều này có nghĩa là dù câu trả lời nghe có vẻ rất thuyết phục, rốt cuộc chúng vẫn dựa trên các mẫu hình trong dữ liệu chứ không phải sự thấu hiểu thực sự (Bockting et al., 2023). Chẳng hạn, khi bạn yêu cầu AI tạo bản tóm tắt nghiên cứu, nó không phân tích phê phán các nguồn như một con người sẽ làm, nó xây dựng câu trả lời dựa trên các cấu trúc ngôn ngữ và xác suất sẵn có. Sự phân biệt này rất quan trọng cần ghi nhớ, vì nó ảnh hưởng đến cả điểm mạnh lẫn hạn chế của công việc học thuật có AI hỗ trợ.

\section{Xây dựng kỹ năng AI}

Xây dựng kỹ năng AI cũng giống như học đi xe đạp. Ban đầu có thể thấy bực bội, nhất là khi bạn không áp dụng đúng kỹ thuật. Như đã đề cập trước đó, sự thất vọng với những kết quả ban đầu có thể dẫn đến việc bỏ cuộc và bỏ lỡ cơ hội tận dụng AI một cách hiệu quả (Mollick \& Mollick, 2023). Cũng như học đi xe đạp đòi hỏi phải hiểu về sự cân bằng, khả năng kiểm soát và chuyển động, việc sử dụng AI tạo sinh đòi hỏi phải thành thạo các prompt (câu lệnh), tinh chỉnh kết quả đầu ra, đồng thời nhận ra điểm mạnh và hạn chế của nó. Mấu chốt là xem đây là một kỹ năng sẽ tiến bộ nhờ luyện tập.

Cuốn sách này cung cấp một nền tảng vững chắc, đặc biệt cho những ai thoải mái với việc tự học. Tuy nhiên, thực hành là điều thiết yếu. Những lời khuyên được nêu ở đây cần được áp dụng, thử nghiệm và điều chỉnh cho phù hợp với nhu cầu cụ thể của bạn. Khi các công cụ AI phát triển nhanh chóng, các ứng dụng đặc thù theo từng lĩnh vực sẽ tiếp tục xuất hiện, và dù chúng tôi không thể bao quát mọi diễn biến mới, các nguyên tắc cốt lõi vẫn không thay đổi. Nếu bạn thường xuyên sử dụng AI tạo sinh và đưa nó vào thói quen làm việc học thuật của mình, bạn sẽ ngày càng thành thạo, cũng như việc đi xe đạp trở thành bản năng thứ hai khi bạn càng luyện tập nhiều (Southworth et al., 2023).

Ngoài việc tự học, việc tìm kiếm thêm các khóa đào tạo và tài nguyên cũng rất có giá trị. Nhiều cơ sở giáo dục đại học hiện nay tổ chức các hội thảo chuyên đề, chương trình đào tạo hoặc khóa học trực tuyến về AI dành cho giảng viên và nhà nghiên cứu (ví dụ: Cornell University, 2023). Các trường đại học ngày càng đưa năng lực AI (AI literacy) vào phát triển chuyên môn, nhận thức được vai trò ngày càng lớn của nó trong nghiên cứu, giảng dạy và quản trị (Carvalho et al., 2022). Bên cạnh đó, còn có vô số tài nguyên công khai, từ các video hướng dẫn trên YouTube đến các khóa học trực tuyến và các buổi thảo luận do chuyên gia dẫn dắt. Như với bất kỳ công nghệ đang phát triển nào, việc luôn cập nhật thông tin và không ngừng trau dồi kỹ năng sẽ giúp bạn sử dụng AI hiệu quả hơn theo thời gian.

Phát triển và duy trì năng lực AI là điều then chốt đối với giới học thuật đang định hướng trong một thế giới ngày càng do AI chi phối. Điều này không chỉ có nghĩa là hiểu cách sử dụng AI tạo sinh mà còn là đánh giá một cách phê phán các kết quả đầu ra của nó, nhận ra những thiên kiến tiềm ẩn, và cập nhật các công cụ mới cũng như các vấn đề đạo đức mới nổi (Bockting et al., 2023). AI không phải là một công nghệ tĩnh tại, mà luôn không ngừng phát triển. Những ai dành thời gian học hỏi, thử nghiệm và tích hợp AI vào quy trình làm việc của mình sẽ ở vị thế tốt nhất để khai thác trọn vẹn tiềm năng của nó. Cũng như năng lực số trở thành một kỹ năng thiết yếu trong giới học thuật cùng với sự trỗi dậy của internet, năng lực AI đang trở nên không thể thiếu không kém (Southworth et al., 2023).

\section{Đánh giá các công cụ AI}

Việc đánh giá các công cụ AI đòi hỏi phải hiểu cả các mô hình AI tạo sinh đa dụng lẫn những ứng dụng chuyên biệt hơn được thiết kế cho công việc học thuật. Dù các mô hình AI đa dụng như ChatGPT, Claude và Gemini tiếp tục được cải thiện, chúng vẫn có những hạn chế, nhất là với các tác vụ đòi hỏi quyền truy cập vào cơ sở dữ liệu độc quyền, phương pháp luận đặc thù của từng lĩnh vực, hoặc kiến thức chuyên ngành (Dwivedi et al., 2023). Khi AI tiến bộ, có thể các công cụ đa dụng này sẽ đủ năng lực hơn để xử lý các tác vụ học thuật chuyên biệt, nhưng hiện tại, các nhà nghiên cứu thường cần tìm hiểu thêm những nền tảng dựa trên AI được thiết kế riêng cho nhu cầu cụ thể của mình (Van Dis et al., 2023).

Các công cụ AI chuyên biệt có thể giúp lấp những khoảng trống mà các mô hình đa dụng chưa đáp ứng được. Chẳng hạn, trong khi ChatGPT có thể hỗ trợ tóm tắt các khái niệm, nó có thể không đủ để thực hiện một tổng quan tài liệu toàn diện, vì không có quyền truy cập trực tiếp vào các cơ sở dữ liệu học thuật. Đây là lúc các nền tảng như Elicit.org phát huy tác dụng, cho phép người dùng tinh chỉnh câu hỏi nghiên cứu, truy xuất các bài báo liên quan và phân tích các phát hiện một cách có hệ thống (OECD, 2023). Một ví dụ khác là Undermind.ai, đóng vai trò trợ lý nghiên cứu chạy bằng AI, quét và tóm tắt khối lượng lớn các bài báo học thuật để cung cấp cho nhà nghiên cứu những hiểu biết sâu sắc hơn. Các công cụ này có thể nâng cao đáng kể hiệu quả, đặc biệt trong những lĩnh vực phụ thuộc nhiều vào phân tích tài liệu quy mô lớn.

Một số công cụ AI thậm chí được thiết kế cho các ngành rất chuyên biệt. Chẳng hạn trong lĩnh vực luật, Harvey AI hỗ trợ các chuyên gia pháp lý phân tích hợp đồng, tóm tắt án lệ và nghiên cứu pháp lý, tinh giản quy trình làm việc theo những cách mà các mô hình AI đa dụng không làm được. Các ứng dụng AI đặc thù theo ngành tương tự cũng đang xuất hiện trong y học, kỹ thuật và khoa học nhân văn, mang lại sự hỗ trợ được thiết kế riêng, phù hợp với nhu cầu nghiên cứu của các cộng đồng học thuật khác nhau (OECD, 2023).

Vì bức tranh về các công cụ AI đang thay đổi nhanh chóng, không có danh sách chuẩn nào về các nền tảng mà mọi nhà nghiên cứu nên dùng. Thay vào đó, các học giả phải đánh giá từng công cụ một cách riêng biệt, xét mức độ hữu ích của nó trong lĩnh vực và quy trình làm việc của chính họ (Bockting et al., 2023). Việc chọn đúng công cụ mang tính chủ quan cao: điều hiệu quả với nhà nghiên cứu này có thể không hiệu quả với nhà nghiên cứu khác. Thử nghiệm các nền tảng khác nhau và đánh giá điểm mạnh, điểm yếu của chúng là điều thiết yếu, nhưng để làm vậy cần có hiểu biết vững chắc về AI tạo sinh. Nếu không nắm vững nền tảng về cách các công cụ này vận hành, rất dễ bỏ sót hạn chế của chúng hoặc sử dụng sai, dẫn đến kết quả dưới mức tối ưu.

Để tận dụng tối đa AI trong công việc học thuật, các nhà nghiên cứu nên cập nhật những diễn biến mới, tham gia các mạng lưới học thuật tập trung vào AI và thử nghiệm các công cụ khác nhau để xác định công cụ nào phù hợp nhất với nhu cầu của mình (Van Noorden \& Perkel, 2023). Cũng như với mọi công nghệ mới nổi, sự thành thạo đến từ thực hành, và những ai chủ động trau dồi năng lực hiểu biết về AI (AI literacy) sẽ có vị thế tốt hơn để khai thác trọn vẹn tiềm năng của nó.

\section{Tích hợp AI vào quy trình làm việc}

Chúng ta đã tìm hiểu cách AI tạo sinh có thể nâng cao hiệu quả viết học thuật, soạn đề xuất tài trợ và xuất bản, nhưng tiện ích của nó vượt xa những nhiệm vụ này. Một khi đã phát triển năng lực hiểu biết về AI, bạn có thể đưa nó vào hầu như mọi khía cạnh trong quy trình làm việc học thuật, giúp tinh giản cả công tác nghiên cứu lẫn trách nhiệm hành chính (Southworth et al., 2023). AI không chỉ là công cụ để viết, nó còn có thể hỗ trợ động não ý tưởng, tự động hóa các tác vụ lặp đi lặp lại, cải thiện giao tiếp và tối ưu hóa việc quản lý thời gian (Bockting et al., 2023).

Một ứng dụng có thể áp dụng ngay là quản lý email và giao tiếp. Giới học thuật liên tục nhận được vô số email, lời mời gửi bài, thư mời hội thảo chuyên đề, cập nhật hành chính, thông báo của tạp chí và thắc mắc của sinh viên. Việc xử lý những email này một cách hiệu quả có thể là một thách thức. AI có thể giúp bằng cách tóm tắt các email dài, cho phép bạn nắm nhanh những điểm chính mà không phải đọc từng chữ (Mollick \& Mollick, 2023). Chẳng hạn, nếu bạn nhận được nhiều thông báo từ nhiều tổ chức khác nhau, AI có thể chỉ trích xuất những thông tin liên quan nhất. Bạn cũng có thể dùng AI để xác định những cơ hội đáng để mình quan tâm bằng cách đưa ra prompt (câu lệnh) như: ``Lời mời gửi bài này có phù hợp với mối quan tâm nghiên cứu của tôi không?'' hoặc ``Hãy tóm tắt thư mời hội nghị này và nêu bật lý do nó có thể phù hợp với tôi.''

Ngoài việc đọc email hiệu quả hơn, AI còn có thể soạn thư trả lời, tiết kiệm thời gian đồng thời bảo đảm sự rõ ràng và nhất quán về giọng điệu. Khi soạn một email, hãy nêu cụ thể điều bạn cần. Chẳng hạn, bạn có thể yêu cầu AI soạn một phản hồi ân cần cho yêu cầu của sinh viên, hoặc một thư trả lời cứng rắn nhưng chuyên nghiệp gửi biên tập viên. AI thậm chí có thể tinh chỉnh giọng điệu, khiến thông điệp lịch sự hơn 10\% hoặc quyết đoán hơn đôi chút. Nếu bạn cần từ chối một lời mời nhưng vẫn muốn giữ mối quan hệ tốt, AI có thể tạo ra một phản hồi vừa thể hiện sự trân trọng vừa mang tính chiến lược. Tương tự, với các email thường nhật, chẳng hạn sắp xếp lịch họp hoặc trả lời các thắc mắc hành chính, bạn chỉ cần dán nội dung email vào AI và yêu cầu nó soạn một thư trả lời ngắn gọn.

AI cũng có thể cải thiện công tác giảng dạy và quản lý khóa học. Nó có thể hỗ trợ thiết kế đề cương môn học, tạo bài giảng và xây dựng thang đánh giá (rubric) (Carvalho et al., 2022). Một ứng dụng đặc biệt có giá trị là biên soạn đề thi. AI có thể giúp tạo câu hỏi trắc nghiệm, câu hỏi trả lời ngắn hoặc đề tài tiểu luận dựa trên tài liệu của khóa học (Pomerlau, 2023). Bạn có thể chỉnh sửa các câu hỏi do AI tạo ra để bảo đảm chúng phù hợp với mục tiêu học tập, mức độ khó và chủ đề của khóa học. Ngoài ra, AI có thể hỗ trợ chấm điểm, đặc biệt với các bài tập có cấu trúc như bài thi trắc nghiệm hoặc các bài đánh giá chuẩn hóa. Với các câu trả lời mở, AI có thể đưa ra phản hồi sơ bộ, nêu bật những điểm chính cần cải thiện, thậm chí gợi ý các thang chấm điểm khả dĩ (Cornell Guidance, 2023). Tuy nhiên, khi dùng AI để chấm điểm, điều thiết yếu là phải xem xét nó đã đánh giá bài làm của sinh viên như thế nào và vì sao. Các gợi ý chấm điểm do AI tạo ra cần được kiểm tra chéo để bảo đảm chúng phù hợp với kỳ vọng học thuật và không đưa vào các thiên kiến (Bockting et al., 2023).

Trong quản lý nghiên cứu, AI có thể hỗ trợ theo dõi và tổ chức tài liệu, giúp giới học thuật cập nhật những diễn biến mới nhất trong lĩnh vực của mình. Các công cụ dựa trên AI có thể tóm tắt các ấn phẩm mới, trích xuất những phát hiện chính và làm nổi bật các xu hướng liên quan (OECD, 2023). Bạn có thể dùng AI để tạo bản tóm lược nghiên cứu bằng cách dán các phần tóm tắt (abstract) và yêu cầu so sánh giữa nhiều bài báo. Nó cũng có thể hỗ trợ cấu trúc hóa các dự án nghiên cứu bằng cách lập lộ trình thời gian, chia nhỏ nhiệm vụ và đưa ra khuyến nghị về các hướng tiếp cận phương pháp luận.

Ngoài nghiên cứu và giảng dạy, AI có thể tinh giản các công việc hành chính. Nhiều học giả phải gánh vác đồng thời nhiều trách nhiệm, từ thẩm định hồ sơ xin tài trợ đến tổ chức hội nghị và quản lý các báo cáo của cơ quan (Southworth et al., 2023). AI có thể hỗ trợ soạn chương trình họp, tóm tắt biên bản cuộc họp, tạo báo cáo, hoặc thậm chí tự động hóa một phần quy trình phản biện đồng cấp (peer review) bằng cách phân tích xu hướng các bài nộp (Conroy, 2023). Theo thời gian, các thiết chế có thể tích hợp AI vào hệ thống nội bộ để nâng cao hiệu quả trên nhiều chức năng học thuật (Van Noorden \& Perkel, 2023).

Chìa khóa để tích hợp AI trọn vẹn vào quy trình làm việc học thuật là thử nghiệm với các ứng dụng khác nhau và đánh giá xem chúng mang lại giá trị lớn nhất ở đâu. Dù AI có thể tự động hóa một số tác vụ, nó hiệu quả nhất khi được dùng như một công cụ cộng tác, công cụ nâng cao hiệu suất mà không thay thế tư duy phản biện và việc ra quyết định. Khi các công cụ AI tiếp tục phát triển, những ai chủ động đưa chúng vào công việc hằng ngày sẽ tìm ra những cách mới để tối ưu hóa năng suất và giải phóng thời gian cho những theo đuổi trí tuệ ở tầm cao hơn (Van Dis et al., 2023).

Xử lý lỗi

AI tạo sinh mang lại những lợi thế đáng kể trong công việc học thuật, nhưng cũng đặt ra một số thách thức đòi hỏi sự chú ý cẩn trọng. Một trong những vấn đề quan trọng nhất là hiện tượng ảo giác (hallucination) (được thảo luận ở Chương 3), khi AI tạo ra thông tin sai lệch hoặc gây hiểu lầm nhưng trông có vẻ đáng tin, trong khi thực chất hoàn toàn bịa đặt (Van Dis et al., 2023). Điều này đặc biệt gây vấn đề trong viết học thuật khi AI bịa ra các trích dẫn, gán sai nguồn hoặc cung cấp dữ liệu không chính xác. Nếu không được kiểm soát, những lỗi như vậy có thể làm suy giảm uy tín của nghiên cứu và dẫn đến sự lan truyền thông tin sai lệch.

Một hạn chế khác là tính chất nông cạn hoặc chưa đầy đủ của phần phân tích do AI tạo ra. Mặc dù AI có thể tóm tắt nội dung, nhưng nó không hiểu trọn vẹn bối cảnh, phương pháp luận hay những sắc thái của các cuộc tranh luận học thuật (Dwivedi et al., 2023). Do đó, các bản tóm tắt do AI tạo ra có thể đơn giản hóa quá mức những lập luận phức tạp hoặc bỏ sót các chi tiết then chốt vốn rất quan trọng đối với thảo luận học thuật. Vấn đề này đặc biệt rõ trong các bài tổng quan tài liệu (literature review), nơi AI có thể ưu tiên những bài báo được trích dẫn thường xuyên hơn là các nghiên cứu mang tính đổi mới hoặc làm thay đổi cả lĩnh vực, qua đó củng cố các thứ bậc học thuật hiện có (Bockting et al., 2023).

AI cũng có thể gây ra sai sót trong các tác vụ định lượng. Khi được yêu cầu thực hiện phép tính, truy xuất dữ liệu thống kê hoặc phân tích xu hướng số liệu, AI có thể tạo ra các con số không chính xác, áp dụng sai công thức hoặc dựa vào các nguồn đã lỗi thời. Nguyên nhân là các mô hình AI tạo sinh không được thiết kế như những công cụ suy luận; chúng dự đoán văn bản dựa trên các mẫu hình thay vì thực hiện các phép toán hay thống kê một cách chính xác (Mollick \& Mollick, 2023). Các học giả dựa vào AI cho nghiên cứu thực nghiệm phải cẩn thận đối chiếu các kết quả số với bộ dữ liệu gốc và các nguồn đáng tin cậy.

Một vấn đề phổ biến khác là sự thiếu chính xác của trích dẫn. Các công cụ AI thường bịa ra tài liệu tham khảo, trình bày chúng như các bài báo tạp chí, sách hoặc kỷ yếu hội nghị có thật (Conroy, 2023). Nếu một tài liệu tham khảo do AI tạo ra trông có vẻ hợp lý nhưng thực tế không tồn tại, nó có thể đánh lạc hướng các nhà nghiên cứu và làm tổn hại liêm chính học thuật. Như chúng ta đã thảo luận ở Chương 3, ngay cả khi AI truy xuất được các trích dẫn có thật, nó vẫn có thể thay đổi các chi tiết như tên tác giả, năm xuất bản hoặc số trang, đòi hỏi phải kiểm chứng thông tin cẩn thận. Để giảm thiểu rủi ro này, các nhà nghiên cứu nên đối chiếu chéo các trích dẫn trong các cơ sở dữ liệu học thuật như Google Scholar, JSTOR hoặc PubMed trước khi đưa chúng vào công trình của mình (Cornell AI Guidelines, 2023).

Lưu ý đến những sai sót tiềm ẩn này là điều thiết yếu trong mọi ứng dụng học thuật của AI. Dù là soạn thảo bài báo nghiên cứu, chuẩn bị tài liệu giảng dạy hay tóm tắt báo cáo, các học giả phải đánh giá một cách phê phán nội dung do AI tạo ra thay vì chấp nhận nó theo giá trị bề mặt. Một thực hành tốt nhất là xác minh các chi tiết quan trọng thông qua nguồn gốc, tham khảo nhiều phản hồi của AI để kiểm tra tính nhất quán, và tinh chỉnh kết quả đầu ra thông qua tương tác lặp đi lặp lại. Bằng cách giữ vai trò chủ động trong quá trình nghiên cứu và viết, các học giả có thể tận dụng lợi ích của AI đồng thời tránh được những cạm bẫy của tự động hóa.

Kết luận

AI tạo sinh đang làm thay đổi công việc học thuật, mang lại những khả năng mới trong nghiên cứu, viết, giảng dạy và quản trị. Mặc dù không thể thay thế tư duy phản biện hay sự nghiêm cẩn học thuật, nó là một công cụ mạnh mẽ, khi được sử dụng hiệu quả, sẽ nâng cao năng suất và tinh gọn quy trình làm việc. Tuy nhiên, những lợi ích của nó đi kèm với các thách thức đòi hỏi phải xử lý cẩn trọng. Năng lực hiểu biết về AI (AI literacy) là then chốt: hiểu cách các công cụ này vận hành, nhận ra giới hạn của chúng và phát triển các chiến lược để sử dụng chúng hiệu quả. AI không phải là giải pháp thần kỳ mà là một nguồn lực đòi hỏi sự tương tác có suy xét. Những học giả thử nghiệm, tinh chỉnh cách tiếp cận và đánh giá một cách phê phán nội dung do AI tạo ra sẽ phát huy tối đa tiềm năng của nó.

Ngoài việc viết, AI hỗ trợ các quy trình học thuật, từ tổng quan tài liệu và chấm điểm đến các tác vụ hành chính và quản lý nghiên cứu. Tuy nhiên, việc sử dụng có trách nhiệm là thiết yếu: bảo đảm tính chính xác, tránh phụ thuộc quá mức và duy trì sự liêm chính học thuật. Các kết quả đầu ra của AI phải được kiểm chứng, vì công nghệ này có thể hiểu sai thông tin, tạo ra các nguồn ảo (hallucination) và gây ra sai sót. Khi AI tiến hóa, vai trò của nó trong học thuật sẽ ngày càng lớn. Những ai chủ động tham gia và trau dồi kỹ năng AI của mình sẽ đứng ở tuyến đầu của cuộc chuyển đổi này. Điều then chốt không phải là chống lại AI mà là tích hợp nó một cách chiến lược, tận dụng điểm mạnh của nó đồng thời giữ vững chiều sâu, tính độc đáo và sự liêm chính vốn định hình học thuật hàn lâm.
